\section{Background Job trong .NET}

\subsection{Background Job giải quyết vấn đề gì?}

Khi người dùng gọi một API, máy chủ nhận yêu cầu, xử lý và trả kết quả. Khoảng thời
gian đó là vòng đời của một HTTP request. Những việc như dọn dữ liệu mỗi đêm, đồng
bộ dữ liệu định kỳ hoặc liên tục đọc hàng đợi không nên phụ thuộc vào request, vì
chúng vẫn phải tiếp tục sau khi máy chủ đã trả kết quả cho người dùng.

Vì không phụ thuộc vào request, Background Job có thể chạy một lần, chạy định kỳ hoặc
chờ dữ liệu mới để xử lý. Vòng đời của công việc gồm thời điểm bắt đầu, quá trình
thực hiện, thời điểm dừng và việc xử lý lỗi. Cách thực hiện bên trong có thể là tuần
tự, bất đồng bộ hoặc song song tùy loại tác vụ.

Không nên chỉ tạo một \texttt{Task} trong request rồi bỏ mặc. Cách làm này thường
được gọi là \textbf{fire-and-forget}: request không chờ và cũng không theo dõi kết
quả. Nếu ứng dụng dừng, công việc có thể chưa hoàn tất; nếu công việc lỗi, lỗi cũng
khó được ghi nhận và xử lý đúng nơi.

\subsection{Các thành phần có sẵn trong .NET}

.NET cung cấp \textbf{Generic Host} (gọi ngắn là \textbf{Host}) để quản lý quá trình
khởi động và dừng ứng dụng. Background Job được đăng ký với Host dưới dạng
\textbf{hosted service} \cite{ms-generic-host,ms-hosted-services}.

\begin{itemize}
  \item \texttt{IHostedService} là một \textbf{interface}, tức hợp đồng quy định
  các phương thức mà hosted service phải có. \texttt{StartAsync} là phương thức Host
  gọi để khởi động service; \texttt{StopAsync} là phương thức Host gọi khi ứng dụng
  dừng.

  \item \texttt{BackgroundService} là lớp cơ sở đã hiện thực
  \texttt{IHostedService}. Lập trình viên kế thừa lớp này và viết
  \texttt{ExecuteAsync}, phương thức chứa công việc được thực hiện trong thời gian
  service hoạt động.

  \item Worker Service là mẫu project tạo sẵn Host và một
  \texttt{BackgroundService}. Mẫu này tạo cấu trúc ban đầu cho một ứng dụng chuyên
  chạy Background Job.
\end{itemize}

\subsection{Vòng đời của hosted service}

Khi ứng dụng khởi động, Host gọi các hosted service đã đăng ký. Khi ứng dụng chuẩn
bị dừng, Host gửi tín hiệu yêu cầu dừng và chờ chúng kết thúc an toàn.

Tín hiệu dừng được truyền bằng \textbf{cancellation token}, tức một đối tượng mang
trạng thái yêu cầu hủy. Token không cưỡng bức thread kết thúc; code trong service
phải kiểm tra hoặc truyền token cho các thao tác bất đồng bộ để chúng chủ động dừng.

\begin{center}
\resizebox{0.84\textwidth}{!}{\begin{tikzpicture}[
  node distance=8mm,
  every node/.style={font=\small,align=center},
  flowbox/.style={draw,rounded corners=2mm,minimum height=11mm,minimum width=25mm,thick},
  arrow/.style={-{Latex[length=2mm]},thick}
]
\node[flowbox,fill=gray!12,draw=gray!65] (host) {Host\\khởi động};
\node[flowbox,fill=blue!15,draw=blue!55,right=of host] (start) {Gọi\\StartAsync};
\node[flowbox,fill=green!15,draw=green!55!black,right=of start] (execute) {Chạy\\ExecuteAsync};
\node[flowbox,fill=orange!18,draw=orange!70!black,right=of execute] (cancel) {Báo hủy qua\\cancellation token};
\node[flowbox,fill=gray!12,draw=gray!65,right=of cancel] (stop) {Chờ\\StopAsync};
\draw[arrow] (host) -- (start);
\draw[arrow] (start) -- (execute);
\draw[arrow] (execute) -- (cancel);
\draw[arrow] (cancel) -- (stop);
\end{tikzpicture}
}
\captionof{figure}{Vòng đời cơ bản của một hosted service}
\label{fig:background-lifecycle}
\end{center}

Luồng trong \figref{fig:background-lifecycle} diễn ra như sau: Host gọi
\texttt{StartAsync} để khởi động service; phần việc chính chạy trong
\texttt{ExecuteAsync}; khi ứng dụng dừng, Host báo hủy rồi gọi
\texttt{StopAsync}. Service cần ngừng nhận việc mới, kết thúc hoặc dừng an toàn việc
đang làm và giải phóng tài nguyên. Quá trình đó gọi là \textbf{graceful shutdown}.

\subsection{Những vấn đề cần quản lý}

Hosted service chỉ cung cấp nơi chạy công việc theo vòng đời ứng dụng. Phần xử lý
vẫn phải xác định rõ:

\begin{itemize}
  \item \textbf{Thời điểm chạy:} chạy một lần khi ứng dụng khởi động, lặp theo
  khoảng thời gian hay liên tục chờ dữ liệu mới.

  \item \textbf{Dừng công việc:} truyền cancellation token xuống các thao tác đang
  chờ để ứng dụng có thể dừng an toàn.

  \item \textbf{Xử lý lỗi:} ghi nhật ký để biết công việc nào thất bại. Chỉ thử lại
  những lỗi tạm thời và tránh để một lần chạy lại tạo dữ liệu trùng.

  \item \textbf{Khả năng mất công việc:} hosted service không tự lưu job. Nếu một
  job bắt buộc phải tiếp tục sau khi ứng dụng khởi động lại, cần lưu nó bên ngoài
  bộ nhớ hoặc dùng một hệ thống quản lý job.
\end{itemize}

Hosted service quản lý Background Job gắn với vòng đời ứng dụng. Hàng đợi bền vững, lịch sử
chạy và lịch thực thi phức tạp thuộc phạm vi của hệ thống quản lý job chuyên dụng.

\subsection{Thay đổi từ .NET Core 3.1 đến .NET 8}

{\small
\begin{longtable}{|>{\raggedright\arraybackslash}p{0.27\textwidth}|
                  >{\raggedright\arraybackslash}p{0.31\textwidth}|
                  >{\raggedright\arraybackslash}p{0.32\textwidth}|}
\caption{Thay đổi của Background Job từ .NET Core 3.1 đến .NET 8}\label{tab:background-version}\\
\hline
\textbf{.NET Core 3.1} & \textbf{Đến .NET 8} & \textbf{Ý nghĩa khi sử dụng} \\
\hline
Lỗi không được xử lý trong \texttt{ExecuteAsync} có thể làm Background Job ngừng nhưng Host vẫn chạy và không ghi nhật ký lỗi mặc định. & Từ .NET 6, lỗi được ghi vào nhật ký và mặc định làm Host dừng. & Lỗi của Background Job dễ được phát hiện hơn; ứng dụng phải quyết định lỗi nào được xử lý, thử lại hoặc để Host dừng. \\
\hline
Thường dùng \texttt{Timer}, bộ hẹn giờ gọi hàm theo mốc thời gian, hoặc \texttt{Task.Delay} để chờ giữa các lần chạy. & \texttt{PeriodicTimer} có từ .NET 6; đây là bộ hẹn giờ cho phép chờ bất đồng bộ đến lần chạy tiếp theo \cite{ms-periodic-timer}. & Vòng lặp định kỳ dễ đọc, dễ nhận cancellation token và tránh hai lượt vô tình chạy chồng nhau. \\
\hline
Code thường đọc trực tiếp \texttt{DateTime} hoặc đồng hồ hệ thống. & .NET 8 có \texttt{TimeProvider}, lớp cung cấp thời gian và bộ hẹn giờ qua một giao diện có thể thay thế. & Khi kiểm thử có thể dùng thời gian giả thay vì phải chờ đồng hồ thật. \\
\hline
\texttt{IHostedService} chỉ có các mốc khởi động và dừng cơ bản. & .NET 8 có \texttt{IHostedLifecycleService}, interface bổ sung các mốc trước và sau khi khởi động hoặc dừng. & Hữu ích khi nhiều service phải phối hợp thứ tự; Background Job thông thường vẫn dùng \texttt{BackgroundService}. \\
\hline
\end{longtable}
}

Thay đổi đáng chú ý nhất đối với code cũ là cách Host xử lý lỗi chưa được bắt trong
\texttt{ExecuteAsync}. Từ .NET 6, lỗi được ghi vào nhật ký và mặc định làm Host dừng; đây là
thay đổi hành vi có thể ảnh hưởng ứng dụng nâng cấp từ .NET Core 3.1
\cite{ms-hosting-exception}. \texttt{TimeProvider} hỗ trợ kiểm thử code phụ thuộc
vào thời gian; \texttt{IHostedLifecycleService} hỗ trợ phối hợp các mốc vòng đời
chi tiết. Cả hai là API được bổ sung trong .NET 8
\cite{ms-time-provider,ms-hosted-lifecycle}.
